\documentclass[11pt,a4paper]{ltjsarticle}
% Common preamble; compile the four documents with LuaLaTeX.
\usepackage{fix-cm}
\usepackage[margin=25mm]{geometry}
\usepackage{luatexja-fontspec}
\setmainjfont{HaranoAjiMincho-Regular.otf}[BoldFont=HaranoAjiMincho-Bold.otf]
\setsansjfont{HaranoAjiGothic-Medium.otf}[BoldFont=HaranoAjiGothic-Bold.otf]
\usepackage{amsmath,amssymb,amsthm,mathtools,booktabs,enumitem}
\usepackage[unicode,hidelinks]{hyperref}
\setlength{\emergencystretch}{2em}
\setlist{nosep}
\newtheorem{theorem}{定理}[section]
\newtheorem{proposition}[theorem]{命題}
\newtheorem{lemma}[theorem]{補題}
\newtheorem{corollary}[theorem]{系}
\theoremstyle{definition}
\newtheorem{definition}[theorem]{定義}
\newtheorem{remark}[theorem]{注意}
\newtheorem{example}[theorem]{例}
\newtheorem{problem}[theorem]{問題}
\newcommand{\C}{\mathbb C}
\newcommand{\Z}{\mathbb Z}
\newcommand{\Q}{\mathbb Q}
\newcommand{\Cat}{\mathcal A}
\newcommand{\DD}{\mathcal D}
\newcommand{\OO}{\mathcal O}
\newcommand{\TT}{\mathcal T}
\newcommand{\PP}{\mathcal P}
\newcommand{\SSer}{\mathbb S}
\newcommand{\Db}{D^{\mathrm b}}
\newcommand{\Hom}{\operatorname{Hom}}
\newcommand{\End}{\operatorname{End}}
\newcommand{\Ext}{\operatorname{Ext}}
\newcommand{\RHom}{\operatorname{RHom}}
\newcommand{\qgr}{\operatorname{qgr}}
\newcommand{\Gr}{\operatorname{Gr}}
\newcommand{\QGr}{\operatorname{QGr}}
\newcommand{\Tor}{\operatorname{Tor}}
\newcommand{\coh}{\operatorname{coh}}
\newcommand{\thick}{\operatorname{thick}}
\newcommand{\gldim}{\operatorname{gldim}}
\newcommand{\GKdim}{\operatorname{GKdim}}
\newcommand{\rk}{\operatorname{rk}}
\newcommand{\id}{\operatorname{id}}
\newcommand{\Hilb}{\operatorname{Hilb}}
\newcommand{\versiondate}{2026年11月17日・v14}

\renewcommand{\versiondate}{2026年11月18日・v18}
\hypersetup{pdftitle={圏論的Fano指数・Serre根・変形：研究方針}}
\title{圏論的Fano指数・Serre根・非可換変形\\研究方針}
\author{}
\date{\versiondate\\再開は\texttt{00\_research\_plan.md}から}
\begin{document}\maketitle
\section{現在の目標と達成基準}
上位目標は，古典的Fano指数を忠実に拡張する圏論的指数と，その根の実現・変形・認識の理論を作ることである．
数値Grothendieck群上の根，heart-exactな自己同値の根，その非可換変形への持上げは相異なる．
現在の主軸は，Sklyanin族による共通Hesse不変二方向の明示的代数化と，
三角形のrank三平行移動格子を持つ無限反Serre商の高次障害を解明することである．
この方針は，数値整除条件だけでは自己同値が得られず，さらに同じ数値根でも変形挙動が異なるために採る．
既知の圏論的偏極・root pairの定義を新設したとはしない．

達成基準は，相対Serreと冪関係まで含む持上げの判定，実際の非可換代数族での実現・非存在，
根付きmoduliの局所構造，および高次表現論・残余圏・安定性・HMSへの実質的な接続である．
元の六公理からの次数9認識は未解決であり，仮定を定義へ移して解決としない．
既存の条件付き認識と全証明は現行論文に保持する．

\section{現在の文書と論理的な分担}
\begin{description}[style=nextline]
\item[\texttt{papers/fano\_root\_deformations.tex}]
現在の中心原稿．HH作用，全Artin次数の根，相対Fano性，次数2の形式的可換性，次数4・Fermat三次の形式族，P2，安定性・entropyを扱う．
全Artin証明は\texttt{papers/sections/root\_artin\_lifting.tex}を読み込む．
\item[\texttt{papers/normalized\_fano\_index.tex}]
古典的指数の回復，normと巡回群障害，surface-like格子の根，del Pezzo指数，三次元の整数値制約を保持する．
\item[\texttt{papers/nc\_dp9\_recognition.tex}]
六公理，有限生成障害，normal境界，Ext消滅，線束復元と認識の全証明を保持する．
\item[\texttt{research/fano\_index\_connections.md}]
34件の公開一次資料の確認範囲，16の関連領域，現在有効な接続と未解決点．外部PDFは収録しない．
\item[\texttt{research/research\_notes.tex}]
認識・数値論に必要な補助証明，形式模型，反例，失敗方法の障害，旧来の文献確認範囲．論文にある証明を重複しない．
\end{description}
数値論文・認識論文・従来ノートは内容を変更していないため文書版はv14のままであり，
現行bundle v18にはそれぞれの現行版だけを収録する．旧bundleを内包しない．

\section{主要な到達点と証明状態}
\subsection{既存の数値・認識結果}
正規化根の次数集合は滑らかなFano多様体の古典指数の約数である．全数値群上のunipotenceと
構造対象の軌道上の条件は異なる．surface-like擬格子上のunipotent根は$-K$の整除性で，
非正規化rank-preserving根は$t_HF$と$N_r(F)H=-K$で判定する．
$r\mid K^2$，正規化なら$r^2\mid K^2$と偶数次数の追加整除を得る．
非正defectでのP2・quadric格子分類はKuznetsovの既知分類に帰着する．

del Pezzo次数1から9で実現可能な最大exact指数は$1,2,3,4,5,6,1,4,3$である．
一般次数4のexact指数は2だが共通の数値格子には四乗根がある．
三次元の整数値Euler制約は，反標準Euler三次性等の追加仮定の下の必要条件である．
詳細な仮定と証明は数値論文の\texttt{thm:classical}，\texttt{thm:twisted-lattice-root}，
\texttt{thm:delpezzo-exact-index-table}，\texttt{prop:degree-four-lifting-gap}，\texttt{thm:threefold-arithmetic}を参照する．

次数9認識は$E_{-1}=0$と$L_3\in\thick(L_0,L_1,L_2)$への有限還元，正次数normal元からの漸近消滅，
次数1・2の整normal境界，一次生成三次境界，可換第三Veronese境界の認識，
滑らかな高Veroneseからのheart復元，線束対象の数値的一意性を保持する．
全数値rank3やEuler格子のunimodularityは一般公理から導いていない．
phantom型の残余を数値類だけで消してはならない．

\subsection{根と変形の結果}
新原稿\texttt{prop:finite-hh}はPerryのSerre自明性から$s_{HH}^r=1$を得る．
\texttt{thm:fixed}はTodaとLowen--Van den Berghを用い，指定したexact関手が持ち上がる一次変形を
$HH^2(X)^{s_{HH}}$に同定する．特殊ファイバーからheart-exact性が継承されることを本文で証明する．
これは標識付き一次変形であり，自己同値群で割ったmoduliの接空間ではない．

一般の候補heartでは\texttt{cor:relative-root}の追加仮定を残すが，tilting付き幾何では
\texttt{prop:tilting-serre}が全flat Artin変形上のexact相対反Serreを構成する．
\texttt{thm:artin-power}は$HH^1=0$の下で関手の持上げが自動的に根となることを示す．
さらに$HH^3=0$を加えた\texttt{thm:formal-roots}では根付き変形が形式的に滑らかで，
小拡大の持上げは$HH^2(X)^s\otimes I$のtorsorとなる．
証明はrelative Serreの全変形上の存在によって再標識を有限次数化し，affine平均化する．
HH上の有限作用だけからcochain作用を仮定する証明ではない．

\texttt{prop:relative-ample}はArtin反標準ample pairとdel Pezzo Hilbert加群を与える．
\texttt{prop:relative-exceptionals}は既存full例外列と根の軌道関係を相対perfect圏に持ち上げる．
$R$が非reducedの場合の全導来圏を絶対smooth properとはしない．

\texttt{cor:geiser-formal}は次数2のGeiser型根を保つ全形式Abel変形が幾何学的な可換変形であることを示す．
\texttt{cor:degree4-formal}は各符号根の不変一次方向を全次数へ延長する（形式次元3または5）．
\texttt{prop:fermat-formal}はFermat三次の三乗根の固定HH次元が2で，可換1方向とPoisson商1方向を持つことを示す．
元の次数4の10-plane配置は一次集合であり，高次の和集合とはまだ同定していない．
同じ接作用でも固定形式部分が異なる例を\texttt{ex:nonlinear-fixed}に保持する．

\subsection{有限反Serre商・二部グラフ・明示代数族}
\texttt{thm:finite-package}は，$\Gamma/\langle\tau\rangle$が有限なら，指定した複数根の同時変形が
形式的に滑らかで接空間が$HH^2(X)^\Gamma$となることを証明する．
前節と同じtilting，$HH^1=HH^3=0$を仮定する．関手同型類の関係を保つがcoherent群作用は構成しない．

\texttt{thm:bipartite-roots}は，次数4の二座標符号根を可換に選べることとグラフの二部性を同値にする．
$D_5$の有理直交基底$v_i=A_i+K/2$で，辺$ij$の全線束解は
$-K/2+a_iv_i+a_jv_j$，$a_i,a_j\in\Z$，$a_i+a_j$奇数である．
可換性は共有頂点係数の一致と同値で，そのparityが二彩色となる．
\texttt{cor:graph-formal}では形式次元は$7-|V_E|$，全五頂点を覆うと全形式変形が可換となる．
奇閉路は有限商を妨げるが，同時変形の非存在を意味しない．

\texttt{prop:quantum-three-roots}は既知のcyclic skew polynomial族に三つの三乗根を明示する．
Nakayamaの恒等性，PBW平坦性，Hesse不変接空間の$xyz$方向を確認する．
この基準例を次数4・三次曲面の代数化と混同しない．

通常の線束根の全一次関手持上げ，P2の固定HH次元10/4，Gepner位相とFano全圏の非存在，
entropy $(0,\dim X)$も現行論文に保持する．独立新規性は未確定である．

\subsection{Hesse代数化と奇閉路の無限反Serre商}
Sklyaninの既知の非退化族$S_{a,b,c}$において，対角・巡回・次数shiftの三つの三乗根は
全fiberで存在する．$[1:-1:0]$でのKodaira--Spencer写像は
$HH^2(\mathbb P^2)^{d,p}=\langle xyz,x^3+y^3+z^3\rangle$への同型である
（\texttt{prop:hesse-two-param}）．既知のSklyanin--superpotential理論に依存し，
二次元Hesse不変方向を実際に代数化する．$HH^1(\mathbb P^2)\ne0$なので
根付きmoduliの形式普遍性は未証明である．

次数4の三角形において$A_0f_{01},A_1f_{12},A_2f_{02}$から生成される群は
反Serre商上で忠実な三次元affine表現を持ち，平行移動格子は
$\{(m_0,m_1,m_2)\in\Z^3:m_0\equiv m_1\equiv m_2\bmod2\}$である
（\texttt{prop:triangle-lattice}）．これはrank三であり，
rank一の無限二面体群の拡大として見積もることはできない．
共通一次不変部分の次元は4だが，有限群平均化は適用できず，二次以降の同時持上げは未証明である．

\section{優先課題と着手方法}
\begin{enumerate}
\item \textbf{Hesse族と根付きmoduli．} Sklyaninの二方向のKodaira--Spencer同型から，
形式的なrooted groupoidの慣性・座標重複を計算する．次数4の符号根またはFermat三次の2方向も固定し，
tilting quiverの積・関係式と可逆核を有限次数で計算する．BWの代数化条件を確認し，
形式存在から一般ファイバーの指数へ飛躍しない．
\item \textbf{奇閉路の無限商．} rank三平行移動格子の三つのPicard tensorが
非可換変形の二次以降に課す条件を調べる．異なる線束解の再標識も二次・三次まで計算し，
奇閉路と共有頂点係数が不一致の根についてPicard commutatorを二次・三次まで調べる．二部グラフの可換な根は既に同時持上げを得たが，全10-planeの高次交叉は未解決である．
\item \textbf{$HH^1\ne0$と巡回coherence．} P2/quadricで自己同値核とnorm補正を調べる．
冪関係の存在，自然同型の選択，Haniharaの巡回不変性を分離する．
\item \textbf{残余圏・HMS・認識．} 相対例外列のmutation回転とGeiser残余Serre関係を比較する．
HH--Serre，AKOの$B+i\omega$，生成threadへ接続する．一般六公理の数値曲面性，defect，normal元なしの消滅，
生成障害は未解決のまま保持する．
\end{enumerate}

\section{検証・文献調査・再現}
新しい直接証明の主要論理を点検し，指標，affine平均化，nilpotent scalar補正，座標配置，Hilbert級数，非線型固定部分を7本のコードで厳密計算する．
有限計算のPASSは圏論的存在や全証明の形式的保証ではない．既存の長い論文の全証明を今回再監査したとはしない．
一次資料は定理と関連証明，定義・主張，abstractのみ，既存確認の継承に分けて文献ノートに記す．
2026年の未監査preprintは比較に留め，新定理の入力にしない．新規性の網羅的確認は未了である．

再現は\texttt{bash scripts/check.sh}と\texttt{research/reproduce.md}のLuaLaTeXコマンドによる．
依存はPython3，SymPy，LuaLaTeX，latexmk，日本語TeXとHarano Ajiフォント．フォントを配布しない．
出力は\texttt{build/}で，bundleには最終PDFだけを収録する．実際の検証・反映・保存状態はREADMEと配布時の報告に記す．

\section{個別方針と区別する恒常的原則・保存許可}
数学的重要性，新規性，理解の深化，応用可能性，正確性，実現可能性を総合して研究価値を最大化する．
方針変更は現時点の理由と元の目標の解決・未解決部分を書き直し，履歴を蓄積しない．
証明済み，条件付き，未完成，予想，計算観察を区別し，仮定・依存結果・証明箇所・新規性の確認範囲を明示する．
誤りがあれば依存主張と本文・要約・方針を訂正し，有用な反例・障害だけをノートに残す．
一次文献を調査し，未確認の内容・書誌情報を推測しない．

原稿は日本語LuaLaTeXで現在の成果を理解しやすい順に構成し，有効な結果と十分な証明を保持する．
研究ノートは現在有効な補助知識を整理し，作業日誌・重複・会話履歴を保存しない．
再開に必要な情報を現行bundleだけで揃え，記法・ファイル・参照・主張状態を整合させる．
bundleには現行文書，必要なソース・PDF・文献情報・コード・入力・証明書・手順だけを収録し，
旧版・旧ZIP・バックアップ・ログ・キャッシュ・中間生成物を除く．容量上限のために数学や検証可能性を削らない．

研究区切りにはユーザーの追加承諾なくLibraryに現行bundleを保存する許可がある．
既存Libraryファイルの削除は明示指示なしに行わない．GitHub \texttt{uedakazushi/nc-dp9}のmainを正本とし，
PR・Actions・force pushを使わず，区切りで通常のfast-forwardを反映する．
生成・計算・組版・GitHub反映・Library保存は独立に確認し，失敗・未実行・再現制約を正確に報告する．
\end{document}
